\section{Model criteria and arithmetic families}\label{sec:models}

\subsection{Numerical corrections on a semistable witness}

The numerical-model criterion of \propref{prop:numerical} is often
applied through a projective dominant witness. In the general
companion formulation \cite{InternalEndpoints}, the witness carries
a rational virtual bundle whose generic degree-$r$ Chern character
is the pulled-back class and whose degree on every vertical
$r$-dimensional subvariety is zero. The native model-Chern detector
then gives zero obstruction upstairs, and dominant detection gives
zero obstruction downstairs. This formulation permits a witness
of positive relative dimension as well as a generically finite one.

We record three ways of producing this numerical correction.
The model calculations and their supporting residue sequences are
given in \cite{InternalModelCriteria}; all model conditions below
refer to supplied projective models.

\begin{proposition}[Scalar total boundary image]\label{prop:scalar-model}
Let $R=K^\circ$ and let $\mathcal X/R$ be integral, regular,
projective reduced strictly semistable, with smooth geometrically
integral generic fiber,
of relative dimension $n\geq3$, with smooth geometrically integral
components $D_i$, smooth closed strata, connected dual graph, and
one relatively ample class $H$. Fix $1\leq r\leq n$ and an
invertible prime $\ell$. Assume that at every vertex $i$ the sum
of Gysin images
\[
 \sum_{j\ne i}\im\bigl(
 H^{2r-2}(E_{ij,\bar k},\Q_\ell(r-1))
       \longrightarrow H^{2r}(D_{i,\bar k},\Q_\ell(r))\bigr)
\]
has dimension one. Then every geometrically homologically trivial
codimension-$r$ cycle on the generic fiber has zero obstruction.
\end{proposition}

\begin{proof}
Choose a Chow extension $\Theta$. The companion specialization
calculation puts its component class in the displayed total image.
Set
\[
 \lambda_i=\deg_k(i_i^!\Theta\,H_i^{n-r}),\qquad
 w_{ij}=\deg_k((H|_{E_{ij}})^{n-1}),
\]
summing disconnected intersections with their native residue degrees.
The specialization trace gives $\sum_i\lambda_i=0$.
The connected weighted graph has positive weights and hence a rational
solution of
\[
 \sum_jw_{ij}(v_j-v_i)=\lambda_i.
\]
The corrected extension
$\Theta_{\rm ad}=\Theta-\sum_i v_i i_{i*}H_i^{r-1}$ has zero
component class: in a one-dimensional total image its degree against
$H_i^{n-r}$ detects that class. Its pure operational representative
is numerically zero on every vertical $r$-cycle. Apply the numerical
model criterion. A constant change of all $v_i$ adds the principal
total-fiber class multiplied by $H^{r-1}$ and hence changes nothing
in Chow.
\end{proof}

One geometric instance is a supplied two-factor hypersurface
degeneration
\begin{equation}\label{eq:two-factor-model}
 AB+\varpi G=0\quad\text{in }\PP^{n+1}_R,
 \qquad\deg A=a,\quad\deg B=b,\quad\deg G=a+b.
\end{equation}
Assume the two special hypersurfaces, their intersection, and
$\Sigma=(A=B=G=0)$ have the smooth transverse geometrically
integral strata specified by these equations, and that the generic
fiber is smooth geometrically integral. Blowing up $(A,\varpi)$
gives components $\Bl_\Sigma(A=0)$ and $(B=0)$ meeting in
$(A=B=0)$. The blowup and restriction computation in
\cite{InternalModelCriteria} verifies the scalar-image condition in
every degree. Thus these local hypersurfaces have $\mathsf V$.

\begin{proposition}[Positive primitive components]\label{prop:positive-model}
Let a smooth projective fourfold have, after a finite extension, a
projective dominant four-dimensional witness with a regular flat
projective model. Write its special fiber as $\sum_i m_iD_i$, with
$m_i>0$, reduced strict normal-crossings support and smooth closed
strata. Suppose a relatively ample $H$ makes
\[
 \deg(xy)\quad\text{positive definite on}\quad
 P_H^2(D_i)=\ker\bigl(H:N^2(D_i)_\Q\to N^3(D_i)_\Q\bigr)
\]
for every component. Then the original fourfold has $\mathsf V$.
\end{proposition}

\begin{proof}
The companion numerical correction leaves a canonical primitive
residual $R_i$ with
$\sum_i m_i\deg(R_i^2)=0$. This residual identity, including the
native constant-field factors, is the primitive-residual theorem of
\cite{InternalModelCriteria}. Positivity and $m_i>0$ give $R_i=0$
for every $i$. The resulting operational lift is numerically
admissible, so codimension two follows by witness descent.
One-cycles and the endpoint degrees complete the fourfold.
\end{proof}

The verified component classes include abelian fourfolds, rank-two
projective bundles over smooth threefolds, rank-three bundles over
smooth surfaces, rank-four bundles over smooth curves, and $\PP^4$.
For the abelian case the positivity input is \cite{Ancona}; the
projective-bundle cases follow from the numerical projective-bundle
formula and divisor Hodge index. If $f:D'\to D_{k'}$ is finite
surjective and $(D',f^*H_{k'})$ has the required positivity, then
$(D,H)$ has it by pullback and the degree formula.

\begin{proposition}[Finite model symmetries]\label{prop:symmetry-model}
For a regular projective model of relative dimension $n\geq4$,
with the preceding multiplicity and stratum conventions, put
\[
 P_H^2=\bigoplus_i\ker\bigl(H^{n-3}:N^2(D_i)_\Q\to
                                    N^{n-1}(D_i)_\Q\bigr),
 \qquad Q(x,y)=\sum_i m_i\deg(x_iy_iH_i^{n-4}).
\]
Let a finite group act on the model preserving the numerical class
of $H$. If $Q$ is nondegenerate and $P_H^2$ is multiplicity-free
as a rational representation, every codimension-two homological
class has zero obstruction. In dimension four, nondegeneracy is
automatic and the conclusion gives $\mathsf V$.
\end{proposition}

\begin{proof}
The image of the canonical residual map is an invariant totally
isotropic submodule of $P_H^2$, by the equivariant residual theorem
in \cite{InternalModelCriteria}. Every rational simple representation
of a finite group is self-dual, by an averaged positive rational
form. Distinct simple summands are therefore orthogonal for $Q$.
Multiplicity-freeness and nondegeneracy imply that each summand is
nondegenerate, so no nonzero invariant totally isotropic submodule
exists. The residual vanishes. Numerical correction and witness
descent give the asserted degree. In dimension four, numerical
duality and divisor Hodge index give the required nondegeneracy;
the endpoints and one-cycles give all remaining degrees.
\end{proof}

The same proof applies to a sum of rational central isotypic
projectors $e$ when $eP_H^2$ has these properties, and gives the
corresponding cycle subspace $e\CH^2_{\Q,\geomhom}$. Such a
subspace pairs globally with the admitted complementary group.

\subsection{Ordinary quadratic reduction}

\begin{theorem}[Companion ordinary-reduction theorem]\label{thm:ordinary}
Let $K/\Q_p$ be finite for any prime $p$, and let $X/K$ be smooth
projective geometrically integral of dimension at least three.
Suppose $X$ has an integral flat projective model whose geometric
special fiber is integral and smooth apart from finitely many
ordinary quadratic points, meaning that their projectivized
geometric tangent cones are smooth quadrics. Then $\mathsf V(X)$ holds. The same
conclusion holds if such a smooth projective dominant witness is
supplied after a finite extension. The theorem reapplies to all
new cycles over every finite valued extension.
\end{theorem}

\begin{proof}
In odd residue characteristic the normal-form and successive
reduced-point blowup calculation of \cite{InternalOrdinary}
produces a regular reduced quadric forest after finite extension.
Its numerical corrections are compatible with the incidence maps.
The all-characteristic quadric-forest theorem in the same source
handles odd-dimensional ordinary quadratic models at residue two.

For even $n\geq4$ at residue two, use the henselian normal form of
\cite[Lemma 4(ii)]{JannsenSaito}. A common finite extension splits
all the finitely many node polynomials. Translating each by an
integral root gives an actual pointed \'etale local equation
\[
 P(x_1,\ldots,x_n)+z^2+\epsilon\varpi^qz=0,
 \qquad q\geq1,\quad\epsilon\in R^\times,
\]
with unimodular polar form for $P$, where $R$ is the auxiliary
field's valuation ring and $\varpi$ its uniformizer. The common \'etale charts and
exact saturated blowup calculations are proved in
\cite[common \'etale neighborhoods and native degenerate-node chains]
{InternalOrdinary}. Each reduced-point blowup decreases $q$ by
one. After exactly $q$ blowups the exceptional chain consists of
$q-1$ copies of
$\PP_E(\mathcal O_E\oplus\mathcal O_E(-1))$ and a terminal
smooth quadric, where
\[
 E=(\overline P(X)+Z^2=0)\subset\PP^n,
 \qquad
 T=(\overline P(X)+Z^2+\overline\epsilon ZS=0)
       \subset\PP^{n+1}.
\]
All double strata are the smooth odd-dimensional quadric $E$;
all multiplicities are one and the intermediate sections are
disjoint. The same quadric-forest numerical criterion now applies.

In every case numerical witness descent uses the complete common
field degree, and the original curved pair gives the form on the
specified support complement. The dominant-witness version follows
from the same detector. After a finite valued extension the ordinary
geometric special points persist, with possibly larger finite
orders; apply the construction anew to every new cycle.
\end{proof}

Over a number field, good reduction or the ordinary alternative at
every finite completion, also on the stated dominant witnesses,
gives the full height pairing. Local auxiliary extensions are
descended before this global construction and need not come from
one fixed global extension.

\subsection{Tate and abelian-type strata}

\begin{theorem}[Companion all-strata criterion]\label{thm:abelian-strata}
Let $Y,T/K$ be smooth projective geometrically integral. Suppose
$Y$ has a supplied regular projective reduced strictly
semistable model whose every connected closed geometric stratum
has a rational Chow motive of abelian type after finite constants.
Assume weight--monodromy for $Y$ in every degree and every
invertible coefficient prime. If $T$ has potentially good projective
reduction, then $T\times Y$ has $\mathsf V$, also for all new
cycles after every finite valued extension.
\end{theorem}

\begin{proof}
The ramified all-strata theorem in \cite{InternalUniformization}
preserves the stated Chow-motive condition on a projective
semistable repair. Its toroidal calculation expresses new strata
through blowups of old boundary strata and projective bundles over
them. For a closed-boundary blowup the property holds in both
directions; this permits the toroidal factorization steps. The
character lattice of a new stratum is primitive, and its calculation
retains the actual line-bundle transition functions.

For a new cycle over $L/K$, pass to a finite extension where $T$
is good and the strata split. On one abelian carrier, numerical
semisimplicity and its comparison prime supply an algebraic
cohomological inner inverse to the boundary incidence. Tensor it
with the diagonal of the coefficient, so the inverse acts on all
mixed K\"unneth summands. Weight--monodromy gives the odd-weight
condition, while the inverse gives numerical-kernel lifting and
actual-span saturation. Apply the companion model criterion and
descend with the full degree. The final curved-pair correction is
on the original support of the newly specified cycle.
\end{proof}

\begin{corollary}[Drinfeld uniformization]\label{cor:drinfeld}
Let $X_a$ be smooth projective algebraizations of
$\Gamma_a\backslash\Omega_K^{d_a}$, where
$\Gamma_a\subset\operatorname{PGL}_{d_a+1}(K)$ is torsion-free,
discrete and cocompact, and $d_a\geq1$. Then
$T\times\prod_aX_a$ has $\mathsf V$ for every potentially good
projective coefficient $T$, also over every finite extension.
\end{corollary}

\begin{proof}
The projective semistable cover of a Drinfeld quotient has
components and closed strata obtained by the successive
rational-linear-subspace blowups of projective space. Their actual
Chow motives are split Tate, by the blowup formula. Staircase
products retain this property. Ito's theorem \cite{Ito} gives
weight--monodromy, and K\"unneth gives it for the product. Apply
\thmref{thm:abelian-strata} and descend the finite cover.
\end{proof}

One may also include curve or surface factors supplied with
projective strict semistable models all of whose closed strata are
of Chow abelian type. The resulting placewise presentations give
full global heights when available at every finite completion;
the factors in those local presentations may depend on the place.

\begin{theorem}[Totally degenerate strata]\label{thm:totally-degenerate}
Suppose a smooth projective $X/K$ has, over a finite extension
$L/K$, a projective dominant generically finite morphism
$Y\to X_L$, with $Y$ smooth projective geometrically integral
of dimension $\dim X$ and carrying a regular projective reduced
strictly semistable model. On every
connected geometric closed stratum $D$, require:
\begin{enumerate}[label=\textup{(\roman*)}]
\item Each integral Chow group is finitely generated. The compatible
hyperplane class $H$ from a projective embedding of the special
fiber satisfies, for $2c\leq\dim D$ and nonzero rational primitive
$x$ with $H^{\dim D-2c+1}x=0$,
\[
 (-1)^c\deg(xH^{\dim D-2c}x)>0.
\]
\item For every $\ell$ different from the residue characteristic,
odd geometric integral cohomology is torsion and
\[
 (\CH^c(D)/\mathrm{tors})\otimes\Z_\ell
 \simeq H^{2c}(D,\Z_\ell(c))/\mathrm{tors}.
\]
\end{enumerate}
Then $\mathsf V(X)$ holds. Smooth projective toric strata satisfy
these hypotheses. A global full height domain follows if such
witnesses, or another full local criterion, are supplied at every
finite place.
\end{theorem}

\begin{proof}
These are conditions (a),(b) of Raskind--Xarles
\cite[Definition 1]{RX}. Define finitely many Chow generators and
components over one finite residue extension and lift it
unramifiedly. The cycle isomorphisms and Poincar\'e duality identify
native numerical groups with even cohomology, compatibly with the
actual incidence correspondences. Thus every numerical-kernel
test lifts to a closed cohomological support test, by the companion
component-detection theorem \cite{InternalDegenerate}.

The cited Raskind--Xarles weight result kills the relevant odd
local Galois groups. More precisely, for
$W=H^{2t-1}(Y_{\bar L},\Q_\ell(t))$ the inertia invariants have
weights at most $-1$, and the inertia-$H^1$ term has weights at
least $1$. Hence $W^{G_L}=H^1(L,W)=0$. The companion period
comparison identifies the rational detector period with the
local Abel--Jacobi cup-product against the closed support test.
It vanishes in $\Q_\ell$, hence is zero as its original rational
intersection number. The finite detector and exact-support
reconstruction give the result. The global assertion applies
\lemref{lem:two-degrees} to these supplied local witnesses.
\end{proof}

\subsection{A common arithmetic construction}

The dense families below use actual local models, together with an
arithmetic pencil argument. We give the latter separately to specify
the field and density quantifiers.

\begin{lemma}[Congruence pencils]\label{lem:congruence-pencils}
Let a fixed integral flat projective ambient model with geometrically
integral fibers and a section system be supplied over a number
ring. Suppose a finite free coefficient module surjects
onto the section system and its required residue spaces. Suppose
the system has a nonempty fixed-member Lefschetz-pencil open, and
that one can prescribe smooth sections at a finite set of places
and an established local smoothing model at distinguished places.
Require the chosen equations to be nonzero on every special
ambient fiber, and the generic smooth-member open to consist of
geometrically integral members.
After a single finite extension realizing these finite residue data,
the congruence-pencil construction gives a Zariski-dense set of
smooth generic members with the prescribed local models, smooth
reduction at the exceptional places, and smooth or ordinary-quadratic
reduction at all remaining places.
\end{lemma}

\begin{proof}
Choose a smooth reference section $H_0$ and include in a finite
set all rational primes below its bad reductions, its zero special
sections, and residue characteristic two. Prepare the desired
residue sections after finite residue extensions. A monic
polynomial with separable irreducible reductions of one common
degree realizes all these extensions in one number field, unramified
at the chosen primes. This field is fixed from now on.

In the free coefficient module impose the finite congruences for a
second section $H_1$, using the square of each native maximal
ideal to retain a prescribed first-order smoothing. For a
distinguished rational prime $p$ unramified in the starting
number field, this is the modulus $p^2$ used in the applications.
Intersect this
affine congruence lattice with the fixed-$H_0$ Lefschetz open. The
intersection is Zariski dense. For a chosen $H_1$, spread its pencil
outside a finite further set of primes. At those additional primes,
choose the parameter to specialize to the already smooth reference
member; at the originally prepared primes choose it to preserve
the chosen residue member. Chinese remainders give infinitely many
parameters with these requirements. Omit the finitely many singular
generic members.

Outside the spreading set, the pencil has smooth or ordinary
quadratic fibers. At every exceptional prime the selected residue
fiber is smooth or is the prescribed established smoothing.
The nonzero residue equations are non-zero-divisors on the integral
ambient fibers. The relative Cartier criterion therefore makes the
resulting divisors flat over the number ring. Generic geometric
integrality follows from the prescribed open. In each application,
nonzero residue equations are checked at the exceptional primes by
the congruences and elsewhere by independence of the two pencil
sections. Finally a polynomial
vanishing on all constructed members vanishes on each permitted
pencil, since that pencil contains infinitely many of them. Varying
$H_1$ in the dense congruence open proves density in the stated
section space. The finite spreading set may depend on $H_1$; the
fixed field does not.
\end{proof}

\subsection{Hypersurfaces and complete intersections}

The following are the geometric conclusions of the local model
calculations in \cite{InternalDense}. Their arithmetic proofs use
\lemref{lem:congruence-pencils}. Each constructed member has full
heights on all new homological classes after every finite extension
of its stated field.

\begin{theorem}[Dense hypersurface families]\label{thm:dense-hypersurfaces}
For every $n\geq3$ and $D\geq3$, there is a Zariski-dense subset
of the full degree-$D$ hypersurface parameter space over $\Q$
consisting of smooth geometrically integral varieties with full
height domains over every number field.

For $n\geq4$, fix positive $d_1,\ldots,d_m$, $m\geq3$,
$D=\sum_i d_i$, and suppose
\[
 b_n=\sum_{|I|=n+1}\prod_{i\in I}d_i-\binom{m-1}{n}>0.
\]
There is such a dense subset with a distinguished product-profile
degeneration of extreme monodromy rank $\rk N^n=b_n$. For linear
factors this gives $D\geq n+2$ and
$\rk N^n=\binom{D-1}{n+1}$.
\end{theorem}

\begin{proof}
For the ordinary family take
\[
 F_0=\sum_{j=0}^{n+1}X_j^D,\qquad
 J_D=\sum_{j=0}^n X_j^{D-1}X_{j+1}+X_{n+1}^D.
\]
At a prime dividing $D$, descending through the partial derivatives
shows $J_D$ is smooth. At every other prime $F_0$ is smooth.
Choose an integral $F_1$ congruent to $J_D$ at primes dividing
$D$ and in the fixed-$F_0$ Lefschetz open. The members
$DF_0+aF_1$, with Chinese-remainder conditions on $a$, have smooth
reduction at the finitely many exceptional primes, including two,
and ordinary reduction elsewhere. Apply \thmref{thm:ordinary}.
Varying $F_1$ proves the asserted full-space density.

For the product profile, prepare transverse integral factors
$P=\prod_i f_i$ and a smoothing section $Q$. Choose
$q>D$ outside their finite bad set and impose
$F_1\equiv P-DF_0+qQ\pmod{q^2}$ as well. For $a\equiv1$
modulo $q^2$, the exact equation at $q$ is $P+qQ_a$ with
$\bar Q_a=\bar Q$. The ordered product-model calculation in
\cite{InternalDense} gives full local vanishing and the displayed
extreme monodromy rank. Its ramified-extension theorem treats all
new cycles after arbitrary finite extension. At the other primes
use the preceding ordinary argument. The complementary-degree
criterion gives the heights.
\end{proof}

\begin{theorem}[Successive complete-intersection sections]\label{thm:dense-ci}
Let $\mathcal Y\subset\PP^N_{\mathcal O_F}$ be a supplied
integral flat projective relative complete intersection of dimension
$m\geq4$, with smooth geometrically integral generic fiber and
geometrically integral special fibers smooth outside finite sets.
For every $d\geq3$, after one fixed finite extension $E/F$ there
is a dense subset of the full degree-$d$ section system on $Y_E$
whose smooth members have full height domains over all finite
extensions of $E$.

Consequently one can construct complete intersections of any
specified multidegree $(d_1,\ldots,d_t)$ with all $d_i\geq3$
and final dimension at least three, over one fixed finite extension
of any prescribed number field. The first equations may be fixed
so that density holds in the final section system. A split quadric
may be used as the initial ambient variety.
\end{theorem}

\begin{proof}
The Koszul resolution supplies the full section space from ambient
polynomials. Smooth residue sections can be chosen to avoid the
finite ambient singular loci, after finite residue preparation.
The degree-$d$ jets for $d\geq3$ give the fixed-member Lefschetz
open. Apply \lemref{lem:congruence-pencils}; the resulting integral
complete-intersection model has smooth or odd-characteristic
ordinary fibers and good fibers at two. Use
\thmref{thm:ordinary}. Its model has the same finite-singular-locus
property, so one can iterate, fixing the earlier equations before
varying the last one.
\end{proof}

There is also a product-profile final-section version on the fixed
ambient of \thmref{thm:dense-ci}, with $m=n+1\geq5$. If the
ambient degree is $\delta$, the final degree is
$D=\sum_{i=1}^s e_i$ with $s\geq3$, $e_i\geq1$, the companion
construction gives a dense final-section family after one fixed
extension, with
\[
 \rk N^n=\delta\sum_{|I|=n+1}\prod_{i\in I}e_i
                                      -\binom{s-1}{n}.
\]
For example degree-five sections of the split quadric fivefold
give a dense family of type-$(2,5)$ fourfolds with $\rk N^4=1$.
The local product theorem and the same congruence-pencil proof
give their complete height domains.

\subsection{Product ambients and positive fourfold hosts}

\begin{theorem}[Dense product-ambient families]\label{thm:dense-products}
For each fixed set of input data below, the constructed dense set
has a full height domain over all finite extensions of one fixed
finite extension of its stated initial number field.
\begin{enumerate}[label=\textup{(\roman*)}]
\item For every $D\geq5$ and every prescribed number field
$F_{\rm base}$, one fixed finite $F/F_{\rm base}$ supports a
dense subset of the full bidegree-$(D,D)$ system on
$(\PP^1\times\PP^4)_F$. A distinguished product model has
\[
 \rk N^4=(D-1)\binom{D-1}{4}>0.
\]
\item Let $\mathcal S/\mathcal O_F$ be a supplied smooth
projective geometrically integral surface model and let $A$ be
relatively very ample. On $\mathcal S\times\PP^3$, put
$L=A\boxtimes\mathcal O(1)$. After one fixed finite extension
$E/F$, the image of degree-five embedding polynomials has a
dense such subset, with
$\rk N^4=10\deg(A^2)-1$. If those polynomials span the full
section space, this is density in $|L^5|$.
\item Let $\mathcal M/\mathcal O_{F_0}$ be a supplied smooth
projective scheme of relative dimension four over the whole number
ring, with geometrically integral fibers. Fix an odd rational prime $p$ unramified
in $F_0$. At each $v\mid p$ and every finite residue extension
$k'/k(v)$ assume
\begin{equation}\label{eq:host-signature}
 b^-(N^2(\mathcal M_{k'})_\R,\deg xy)
                  =\rho(\mathcal M_{k'})-1.
\end{equation}
Choose $H$ sufficiently positive for relative very ampleness,
normal generation and section base change. Put
$L_0=H\boxtimes\mathcal O_{\PP^1}(1)$ and $c=\deg H^4$.
One fixed finite $F/F_0$, unramified above $p$, supports a dense
subset of the full $|L_0^5|_F$ with
\[
 \rk N^4=4(5c-4)>0,\qquad
 b_1(X)=b_1(M),\quad b_2(X)=b_2(M)+1.
\]
\end{enumerate}
\end{theorem}

\begin{proof}
For (i), use $D$ transverse $\mathcal O(1,1)$ factors at a
distinguished prime. The ambient double-stratum comparison retains
both hyperplane classes. The ordered model and its uniform
finite-extension comparison, proved in \cite{InternalDense}, give
full local vanishing and
$5\binom D5-\binom{D-1}4=(D-1)\binom{D-1}4$.

For (ii), the local surface-ambient comparison uses the actual
blowups of the double strata. Its extreme incidence complex gives
$10\deg(A^2)-1$. The free module of embedding polynomials supplies
the coefficient space in \lemref{lem:congruence-pencils}, proving
density in its image. The stated spanning hypothesis identifies
that image with the full section system.

For (iii), the positive-host suspension calculation in
\cite{InternalPositiveHosts} applies to five divisors on $M$
and five distinct point factors on $\PP^1$. Residue Bertini
prepares smooth geometrically integral intersections of one,
two and three divisors, split finite-\'etale fourfold intersections
of degree $c$, and empty fivefold intersection. The smoothing
section cuts each positive-dimensional stratum transversely,
with geometrically integral positive-dimensional cuts and
finite-\'etale zero-dimensional cuts, and avoids the
zero-dimensional strata. Condition \eqref{eq:host-signature} kills
the primitive numerical residual. Its uniform form over finite
residue extensions gives the all-new-cycle ramified conclusion.
The incidence complex has $\dim H_3=5c-4$ before joining the
five point vertices, hence $\rk N^4=4(5c-4)$ afterward.
Weak Lefschetz gives the two Betti numbers. Relative Serre
vanishing and a sufficiently high Veronese supply the full
coefficient module and the required jets. Apply
\lemref{lem:congruence-pencils} to all three cases, use ordinary
reduction at nondistinguished bad primes, and then
\lemref{lem:two-degrees}.
\end{proof}

The signature hypothesis in (iii) is supplied by everywhere-good
abelian fourfolds, by the projective-bundle host types listed after
\propref{prop:positive-model}, and by $\PP^4$. For $\PP^4$,
$H=\mathcal O(1)$ and $p=5$, the result gives bidegree-$(5,5)$
fourfolds with $\rk N^4=4$, $b_1=0$, $b_2=2$.

\subsection{Explicit elliptic-power families}

Let $E_0/\Q$ be $y^2=x^3-x$ and put
$F_{\rm tor}=\Q(E_0[3])$, adjoining all geometric $3$-torsion.
The elliptic curve has an everywhere-good model $\mathcal E$
over this field. At two, potential good reduction follows from
complex multiplication; after adjoining $3$-torsion its finite
inertia image lies in the torsion-free principal congruence
subgroup of $\operatorname{GL}_2(\Z_3)$, so it is trivial. The
curve is already good at the other primes. This field is
unramified above five.

Put $A=\mathcal O_E(3[0])\boxtimes\mathcal O_E(3[0])$ on $E^2$,
and $H=\boxtimes_{i=1}^4\mathcal O_E(3[0])$ on $E^4$.
The plane-cubic embeddings and the Segre construction prove full
section spanning. The following dense subsets of full linear
systems are constructed over fixed finite extensions of
$F_{\rm tor}$ unramified above five, and have full heights after
every further finite extension:
\begin{center}
\begin{tabular}{llll}
\toprule
Ambient and system & Distinguished profile & Monodromy & $(b_1,b_2)$\\
\midrule
$E^2\times\PP^3$, $|(A\boxtimes\mathcal O(1))^5|$
 & five section factors & $\rk N^4=179$ & $(4,7)$\\
same system & curve and five hyperplanes & $N^4=0$ & $(4,7)$\\
$E^4\times\PP^1$, $|(H\boxtimes\mathcal O(1))^5|$
 & two curves and five points & $\rk N^2=802816$, $N^3=0$ & $(8,29)$\\
same system & five divisors and five points & $\rk N^4=38864$ & $(8,29)$\\
\bottomrule
\end{tabular}
\end{center}
The first row has $\deg A^2=18$. The second uses the mixed
surface--hyperplane local theorem of \cite{InternalMixedFamilies};
its weight-two contribution contains a nonzero eight-dimensional
summand. In the third row the double-stratum model is the actual
blowup of the curve product along the smoothing cuts. The local
signature condition holds for the abelian fourfold product, and
the exact rank formula is
\[
 \rk N^2=(2g(C)-b_1(S))(2g(D)-b_1(T))(d-1),\qquad N^3=0.
\]
Here $g(C)=g(D)=226$, $b_1(S)=b_1(T)=4$, and $d=5$.
The final row is \thmref{thm:dense-products}(iii), with
$c=4!\,3^4=1944$. The arithmetic argument is again
\lemref{lem:congruence-pencils}, now with the full Segre coefficient
system. The different monodromy profiles retain the distinctions
among these constructed subsets even when the ambient system is
the same.
